\documentclass[11pt]{article}
\usepackage[a4paper,margin=1in]{geometry}
\usepackage{amsmath,amssymb}
\usepackage{mathtools}

\title{RaceShowdown Model Notes}
\author{}
\date{}

\begin{document}
\maketitle

\section*{Problem Setup}

In this project, a race is determined by two inputs:
\begin{itemize}
    \item $n \geq 0$: the number of hot dogs both athletes eat,
    \item $d \geq 0$: the running distance in meters after eating.
\end{itemize}

The two athletes are Chestnut and Bolt.
The goal is to determine, for each pair $(d,n)$, which athlete is favored to win.

\medskip

There are two versions of the model in the codebase:
\begin{itemize}
    \item a \textbf{deterministic model}, which compares mean total race times,
    \item a \textbf{probabilistic model}, which adds randomness and estimates win probabilities.
\end{itemize}

\section*{Deterministic Model}

For athlete $i$, define the following parameters:
\begin{itemize}
    \item $v_i > 0$: baseline running speed,
    \item $k_i \geq 0$: running fatigue coefficient,
    \item $a_i > 0$: baseline time for a hot dog,
    \item $b_i \geq 0$: extra eating time per hot dog,
    \item $c_i \geq 0$: linear post-eating running penalty,
    \item $q_i \geq 0$: quadratic post-eating running penalty.
\end{itemize}

\subsection*{Eating Time}

If athlete $i$ eats $n$ hot dogs, then the deterministic eating time is
\[
E_i(n)=\sum_{m=1}^{n}(a_i+b_i m).
\]

This can be simplified to
\[
E_i(n)=n a_i + b_i \frac{n(n+1)}{2}.
\]

In the code, this corresponds to the cumulative eating-time calculation.

\subsection*{Post-Eating Running Penalty}

After eating $n$ hot dogs, athlete $i$ runs with a multiplicative penalty
\[
M_i(n)=1+c_i n+q_i n^2.
\]

This term captures the idea that running gets harder after eating more hot dogs.

\subsection*{Running Time}

Before accounting for eating, the baseline running time over distance $d$ is
\[
R_i^{\text{base}}(d)=\frac{d}{v_i}(1+k_i d).
\]

The factor $(1+k_i d)$ models fatigue with distance.

After including the post-eating penalty, the deterministic running time becomes
\[
R_i(d,n)=R_i^{\text{base}}(d)\,M_i(n)
=\frac{d}{v_i}(1+k_i d)(1+c_i n+q_i n^2).
\]

\subsection*{Total Deterministic Time}

The deterministic total race time for athlete $i$ is
\[
T_i(d,n)=E_i(n)+R_i(d,n).
\]

For this project, the deterministic comparison is based on the difference
\[
\Delta(d,n)=T_C(d,n)-T_B(d,n),
\]
where $C$ denotes Chestnut and $B$ denotes Bolt.

The sign of $\Delta(d,n)$ determines the winner:
\begin{itemize}
    \item $\Delta(d,n)<0$: Chestnut is faster,
    \item $\Delta(d,n)>0$: Bolt is faster,
    \item $\Delta(d,n)=0$: deterministic tie.
\end{itemize}

\subsection*{Expanded Deterministic Formula}

Expanding the total-time difference gives
\[
\Delta(d,n)=\bigl(E_C(n)-E_B(n)\bigr)
+ d\left(\frac{M_C(n)}{v_C}-\frac{M_B(n)}{v_B}\right)
+ d^2\left(\frac{k_C M_C(n)}{v_C}-\frac{k_B M_B(n)}{v_B}\right).
\]

This equation is useful because it separates the race into three effects:
\begin{enumerate}
    \item \textbf{Eating advantage:}
    \[
    E_C(n)-E_B(n).
    \]
    This measures who is faster at eating $n$ hot dogs.

    \item \textbf{Linear running effect in distance:}
    \[
    d\left(\frac{M_C(n)}{v_C}-\frac{M_B(n)}{v_B}\right).
    \]
    This captures how the baseline running-speed difference interacts with the post-eating penalty.

    \item \textbf{Quadratic fatigue effect in distance:}
    \[
    d^2\left(\frac{k_C M_C(n)}{v_C}-\frac{k_B M_B(n)}{v_B}\right).
    \]
    This captures the extra long-distance fatigue effect, which becomes more important at large $d$.
\end{enumerate}

This decomposition explains why the winner can change across the $(d,n)$ plane. One athlete may have an advantage in eating, while the other may recover that advantage through running speed or long-distance fatigue behavior.

\section*{Probabilistic Model}

The probabilistic model keeps the same overall structure, but now total race time is random rather than fixed. In the code, randomness is introduced through:
\begin{itemize}
    \item eating-time noise,
    \item running-speed noise,
    \item fatigue noise,
    \item post-eating penalty noise.
\end{itemize}

So instead of deterministic total times $T_C(d,n)$ and $T_B(d,n)$, we consider random total times, which by abuse of notation we still write as $T_C(d,n)$ and $T_B(d,n)$.

The main object of interest is
\[
p(d,n)=\Pr\bigl(T_C(d,n)<T_B(d,n)\bigr).
\]

This is the probability that Chestnut beats Bolt at the point $(d,n)$.

\subsection*{Interpretation of $p(d,n)$}

\begin{itemize}
    \item If $p(d,n)>0.5$, then Chestnut is more likely to win.
    \item If $p(d,n)<0.5$, then Bolt is more likely to win.
    \item If $p(d,n)=0.5$, then both athletes are equally likely to win.
\end{itemize}

In the code, this probability is estimated by Monte Carlo simulation:
\begin{itemize}
    \item sample many race times for Chestnut,
    \item sample many race times for Bolt,
    \item count the fraction of samples in which Chestnut's time is smaller.
\end{itemize}

\section*{What the Two Boundaries Mean}

There are two different boundaries in the project:

\subsection*{Deterministic Boundary}

The deterministic boundary is defined by
\[
\Delta(d,n)=0.
\]

This is the set of points where the athletes have the same \emph{mean-model} race time.

\subsection*{Probabilistic Boundary}

The probabilistic boundary is defined by
\[
p(d,n)=0.5.
\]

This is the set of points where the athletes are equally likely to win once randomness is included.

\medskip

These boundaries do not have to coincide. Two athletes can have the same mean time but different variances, or one athlete can have a slightly worse mean but still win often because of the shape of the distribution.

\section*{Why Non-Monotonic Behavior Can Appear}

It is tempting to expect that if the number of hot dogs is fixed, then increasing the distance should move the outcome in only one direction. However, the deterministic difference
\[
\Delta(d,n)=A(n)+B(n)d+C(n)d^2
\]
is generally a quadratic function of $d$, where
\[
A(n)=E_C(n)-E_B(n),
\]
\[
B(n)=\frac{M_C(n)}{v_C}-\frac{M_B(n)}{v_B},
\]
\[
C(n)=\frac{k_C M_C(n)}{v_C}-\frac{k_B M_B(n)}{v_B}.
\]

Since a quadratic need not be monotone, the deterministic tie boundary is not guaranteed to be monotone either.

The same is true, a fortiori, for the probabilistic boundary, because it depends on full random race-time distributions rather than only on mean values.

\section*{Summary}

The project studies a two-parameter race model in the variables $(d,n)$:
\begin{itemize}
    \item the deterministic model compares mean total times through $\Delta(d,n)$,
    \item the probabilistic model compares win likelihood through $p(d,n)$.
\end{itemize}

The deterministic formula
\[
\Delta(d,n)=\bigl(E_C(n)-E_B(n)\bigr)
+ d\left(\frac{M_C(n)}{v_C}-\frac{M_B(n)}{v_B}\right)
+ d^2\left(\frac{k_C M_C(n)}{v_C}-\frac{k_B M_B(n)}{v_B}\right)
\]
shows how eating, running speed, and fatigue interact.

The probabilistic quantity
\[
p(d,n)=\Pr\bigl(T_C(d,n)<T_B(d,n)\bigr)
\]
is the stochastic analogue: it measures how often Chestnut wins once randomness is added to the race.

\end{document}
